\renewcommand{\chaptertitle}{Combinators for Fun and Profit}

%  Make first page footer:
\fancypagestyle{firststyle}{%
\fancyhf{}% Clear header/footer
\fancyhead{}
\fancyfoot[L]{{\footnotesize
              %% We toggle between these:
              % Manuscript submitted for review.\\
              \authorname{}~\authorpatp{}.  ``\chaptertitle{}''.  {\it The Nock Book} (2025):  10–$x$.
              }}
}
%  Arrange subsequent pages:
\fancyhf{}
\fancyhead[LE]{{\urbitfont The Nock Book}}
\fancyhead[RO]{\chaptertitle{}}
\fancyfoot[LE,RO]{\thepage}

\chapter{\chaptertitle}

\thispagestyle{firststyle}

\begin{abstract}

\end{abstract}

The first section of this chapter will introduce the basic \textct{SKI} and \textct{BCKW} combinator calculi and examine how they can represent each other.  The second section will introduce Nock's opcodes as combinators, and show how they can be used to express various combinator idioms.  The third section will show how to build practical affordances for common programming tasks, such as conditionals, recursion, subprocedure invocation, and data structures.  The final section will discuss the relationship between combinators and virtualization, and how Nock can be used to build a virtual machine even in a crash-only context.

\section{Combinators in Computing}

The pioneers of computing elaborated three complementary visions:

\begin{enumerate}
  \item  The Turing machine, a simple abstract machine that can compute anything computable.
  \item  The lambda calculus, a pure functional language based on substitution and abstraction.
  \item  Combinatory logic, a variable-free functional language based on combinators.
\end{enumerate}

Logicians early demonstrated the equivalence of these approaches, but particular formulations are still preferred for particular tasks.  The Turing machine is the most intuitive model of computation, and is often used to prove computability results.  The lambda calculus is the basis of many functional programming languages, such as Lisp, Scheme, and Haskell.  Combinatory logic is less widely used in practice, but has a long history in mathematical logic and type theory.  It is also the basis of Nock.

Before we get too far out over our skis, let's talk about what a combinator is.  At its simplest, a combinator is a function that takes one or more arguments and returns a result, without referring to any variables outside its own arguments.  In practice, we'll use different combinators to build up more complex functions, and we'll see informally how they can be combined to express any computable function.  We can even build combinators out of other combinators.  The objective of combinatory logic is to express all computation in terms of combinators alone, without the need for variables or named functions.

This first section will introduce some elemental and compound combinators and make the case that they circumscribe ``computing'' as a concept.

\subsection{The \textct{SKI} calculus}

Computable functions can be expressed using minimal sets of combinators.  Many logicians today prefer to work with Moses Schönfinkel and Haskell Curry's \textct{SKI} combinator calculus, which has three combinators:


The identity combinator \textct{I} is conventionally written \textct{Ix = x}.  This simply means that it returns its single argument unchanged.  There is little to say of this combinator at this point, but logicians have found that it frequently simplifies expressions.

\begin{lstlisting}[style=listingcomb]
I x
x
\end{lstlisting}

The constant combinator \textct{K} is conventionally written \textct{Kxy = x}.  This means that it takes two arguments and returns the first, ignoring the second.  It is useful for building functions that always return a fixed value.

\subsubsection{Questions}

\begin{enumerate}
  \item  What is the result of \textct{I x}?
  \item  What is the result of \textct{I I x}?
\end{enumerate}

\begin{lstlisting}[style=listingcomb]
K x y
x
\end{lstlisting}

The substitution combinator \textct{S} is conventionally written \textct{Sxyz = xz(yz)}.  This means that it takes three arguments: a function \textct{x}, and two values \textct{y} and \textct{z}.  It applies the function \textct{x} to the value \textct{z}, and also applies the function \textct{y} to the value \textct{z}, then applies the result of \textct{x z} to the result of \textct{y z}.  This combinator is more complex than the others, but it is very powerful, as it allows us to express function application and composition.

\begin{lstlisting}[style=listingcomb]
S x y z
x z (y z)
\end{lstlisting}

The \textct{SKI} system is Turing complete \citep{blah}, meaning that any computable function can be expressed using only these three combinators.  It is also equivalent to the lambda calculus, meaning that any lambda expression can be translated into an equi\-valent \textct{SKI} expression and vice versa.  (cite church-turing thesis stuff)

Combinators can be combined to form more complex expressions.  For example, the combinator \textct{S K K} is equivalent to the identity combinator \textct{I}:

\begin{lstlisting}[style=listingcomb]
S K K x
K x (K x)
x
\end{lstlisting}

% https://en.wikipedia.org/wiki/Magnum_opus_(alchemy)

% https://en.wikipedia.org/wiki/B,_C,_K,_W_system
% https://dkeenan.com/Lambda/index.htm

% https://doisinkidney.com/posts/2020-10-17-ski.html
% https://www.youtube.com/watch?v=GawiQQCn3bk

\subsection{The \textct{BCKW} calculus}

The \textct{BCKW} combinator calculus is an alternative to \textct{SKI}, with a different set of primitive combinators.  It has four combinators:

The composition combinator \textct{B} is conventionally written \textct{Bxyz = x(yz)}.  This means that it takes three arguments: a function \textct{x}, and two values \textct{y} and \textct{z}.  It applies the function \textct{y} to the value \textct{z}, then applies the function \textct{x} to the result.

\begin{lstlisting}[style=listingcomb]
B x y z
x (y z)
\end{lstlisting}

% The constant combinator \textct{C} is conventionally written \textct{Cxyz = xzy}.  This means that it takes three arguments: a function \textct{x}, and two values \textct{y} and \textct{z}.  It applies the function \textct{x} to the values \textct{z} and \textct{y}, effectively reversing the order of the last two arguments.

% The weakening combinator \textct{W} is conventionally written \textct{Wxy = xyy}.  This means that it takes two arguments: a function \textct{x} and a value \textct{y}.  It applies the function \textct{x} to the value \textct{y} twice.

% The duplication combinator \textct{U} is conventionally written \textct{Uxyz = xzyz}.  This means that it takes three arguments: a function \textct{x}, and two values \textct{y} and \textct{z}.  It applies the function \textct{x} to the values \textct{z}, \textct{y}, and \textct{z} again.  This combinator is less common, but can be useful in certain contexts.

The \textsc{SKI} and \textsc{BCKW} systems are equivalent, meaning that any expression in one system can be translated into an equivalent expression in the other system.  For example, the identity combinator \textct{I} can be expressed in \textsc{BCKW} as:

\begin{lstlisting}[style=listingcomb]
I = W K

\end{lstlisting}


\textct{U}

\textct{Y}

\begin{lstlisting}[style=listingcomb]
Y = S (K (S I I)) (S (S (K S) K) (K (S I I)))
\end{lstlisting}

\textct{MABT} or whatever it is

Chris Barker
The iota combinator \textct{ι} can produce all the others, but in convoluted ways.

% https://stackoverflow.com/questions/7368507/expressing-y-in-term-of-ski-combinators-in-javascript

Combinator calculi are generally considered too minimalist for practical programming; they shine, however, in mathematical proofs.  This means that we can use them to define a standard of behavior for an algorithm, but implement the logic in a compliant form using whatever computational tools we like.  Thus it is better to think of Nock not as a bare programming language, but as a standard of computation which can under many circumstances be used as the computation itself.\footnote{Consideration of when to do so is largely driven by performance bottlenecks.}

\section{Nock Opcodes as Combinator}

The Nock specification is produced in full in Listing~\ref{lst:nock-4k}.

\begin{lstlisting}[style=listingcode,
                   caption={Nock 4K},
                   label={lst:nock-4k}]
Nock 4K

A noun is an atom or a cell. An atom is a natural number. A cell is an ordered pair of nouns.

Reduce by the first matching pattern; variables match any noun.

nock(a)             *a
[a b c]             [a [b c]]

?[a b]              0
?a                  1
+[a b]              +[a b]
+a                  1 + a
=[a a]              0
=[a b]              1

/[1 a]              a
/[2 a b]            a
/[3 a b]            b
/[(a + a) b]        /[2 /[a b]]
/[(a + a + 1) b]    /[3 /[a b]]
/a                  /a

#[1 a b]            a
#[(a + a) b c]      #[a [b /[(a + a + 1) c]] c]
#[(a + a + 1) b c]  #[a [/[(a + a) c] b] c]
#a                  #a

*[a [b c] d]        [*[a b c] *[a d]]

*[a 0 b]            /[b a]
*[a 1 b]            b
*[a 2 b c]          *[*[a b] *[a c]]
*[a 3 b]            ?*[a b]
*[a 4 b]            +*[a b]
*[a 5 b c]          =[*[a b] *[a c]]

*[a 6 b c d]        *[a *[[c d] 0 *[[2 3] 0 *[a 4 4 b]]]]
*[a 7 b c]          *[*[a b] c]
*[a 8 b c]          *[[*[a b] a] c]
*[a 9 b c]          *[*[a c] 2 [0 1] 0 b]
*[a 10 [b c] d]     #[b *[a c] *[a d]]

*[a 11 [b c] d]     *[[*[a c] *[a d]] 0 3]
*[a 11 b c]         *[a c]

*a                  *a
\end{lstlisting}

\subsection{\textct{I} Identity = Opcode 0}

Nock's binary tree addressing is fully general and yields the identity combinator \textct{I} as a special case.  (The cost, of course, is dragging in some machinery for basic arithmetic, as we will see.)

Conceptually, we could write this in a hybrid form between combinators and Nock thus:

\begin{lstlisting}[style=listingcode]
  *[a I b] == [0 1]
\end{lstlisting}

\subsection{\textct{K} Constant = Opcode 1}

\subsection{\textct{S} Substitution = Opcode 2}

\subsection{The Axiomatic Operators}

Because of Nock's cellular nature as a binary tree, we need the ability to access and manipulate nouns.

opcode 3, 4, 5, 10


\section{Idioms \& Design Patterns}

Several of the combinators we examined above are not present in Nock \emph{simpliciter}, but we can build them out of the other pieces.  There are also common idioms which are baked in as ``compound opcodes'' (such as conditional branching, opcode 6).

\subsection{\textct{C} Reversal}

C = S (S (K (S (K S) K)) S) (K K)

In Nock, the reversal combinator \textct{C} is almost free due to the binary tree nature of cells.

\begin{lstlisting}[style=listingcode]
  *[a C b] == *[7 [0 3] [[0 3] [0 2]]]
  *[a C b] == *[[0 7] [0 6]]
\end{lstlisting}


Recursion we get for free with opcode 0.

$\mu$-recursive functions


\section{Virtualization}

Nock is a crash-only pure language, which means at the base level it can only succeed and yield a new noun:  no side effects, no exceptions, no state.  Like a castaway on a desert island, we must build everything we need from raw materials at hand.

Nock has one trick built in and a couple of design patterns to help us get there:

\begin{enumerate}
  \item  Opcode 11 lets us send a hint to the runtime, which can be used to trigger side effects.  The Nock expression formally doesn't know anything about these, but we can use them to print results, read input, and so forth.
  \item  The runtime can maintain state outside of Nock and supply it as a subject to subsequent Nock expressions.  We will do a lot more with this later, but we want to highlight it as a key affordance now.
  \item  The virtualization pattern means that we can run Nock in Nock.
\end{enumerate}

Opcode 11 can encode anything we want:  the result of a hint is used by the runtime, and so anything the runtime is willing to work with can be produced.


The magic formula at the beating heart of a practical Nock operating system is simply:

\begin{lstlisting}[style=listingcode]
*([state -.events] [9 2 10 [6 0 3] 0 2])
\end{lstlisting}

\noindent
That is, apply the Nock formula against a subject pair consisting of the current state and the head of the list of events to process.
